\documentclass[11pt]{article}

% -------------------------------------------------
% Page Layout
% -------------------------------------------------
\usepackage[margin=1in]{geometry}

% -------------------------------------------------
% Fonts & Encoding
% -------------------------------------------------
\usepackage[T1]{fontenc}
\usepackage[utf8]{inputenc}
\usepackage{lmodern}

% -------------------------------------------------
% Mathematics
% -------------------------------------------------
\usepackage{amsmath}
\usepackage{amssymb}
\usepackage{amsfonts}
\usepackage{amsthm}
\usepackage{mathtools}
\usepackage{bm}
\usepackage{mathrsfs}

% -------------------------------------------------
% Figures & Tables
% -------------------------------------------------
\usepackage{graphicx}
\usepackage{float}
\usepackage{booktabs}
\usepackage{array}
\usepackage{multirow}
\usepackage{tabularx}

% -------------------------------------------------
% Lists
% -------------------------------------------------
\usepackage{enumitem}

% -------------------------------------------------
% Colors & Hyperlinks
% -------------------------------------------------
\usepackage{xcolor}
\usepackage[
    colorlinks=true,
    linkcolor=blue,
    citecolor=blue,
    urlcolor=blue
]{hyperref}

% -------------------------------------------------
% Page Style
% -------------------------------------------------
\usepackage{fancyhdr}
\pagestyle{fancy}
\fancyhf{}
\rhead{\thepage}
\renewcommand{\headrulewidth}{0.4pt}

% -------------------------------------------------
% Spacing
% -------------------------------------------------
\usepackage{setspace}
\onehalfspacing

% -------------------------------------------------
% Useful Commands
% -------------------------------------------------
\newcommand{\R}{\mathbb{R}}
\newcommand{\C}{\mathbb{C}}
\newcommand{\N}{\mathbb{N}}
\newcommand{\Z}{\mathbb{Z}}
\newcommand{\Q}{\mathbb{Q}}

\DeclareMathOperator{\rank}{rank}
\DeclareMathOperator{\spanof}{span}
\DeclareMathOperator{\nullity}{null}
\DeclareMathOperator{\trace}{tr}
\DeclareMathOperator{\diag}{diag}
\DeclareMathOperator{\Var}{Var}
\DeclareMathOperator{\Cov}{Cov}
\DeclareMathOperator{\E}{\mathbb{E}}
\DeclareMathOperator{\Prob}{\mathbb{P}}

% -------------------------------------------------
% Theorem Environments
% -------------------------------------------------
\theoremstyle{definition}
\newtheorem{definition}{Definition}[section]
\newtheorem{example}[definition]{Example}
\newtheorem{exercise}[definition]{Exercise}

\theoremstyle{plain}
\newtheorem{theorem}[definition]{Theorem}
\newtheorem{lemma}[definition]{Lemma}
\newtheorem{proposition}[definition]{Proposition}
\newtheorem{corollary}[definition]{Corollary}

\theoremstyle{remark}
\newtheorem{remark}[definition]{Remark}

% -------------------------------------------------
% Title Information
% -------------------------------------------------
\title{\textbf{Introductory Exam Notes}}
\author{Jonathan Ma}
\date{\today}

% =================================================
\begin{document}

\maketitle

\tableofcontents

\newpage
%================================================
\part{Linear Algebra}

\section{Vectors \& Analytic Geometry of Space}

\subsection{Coordinate Systems in $n$-space}

\subsection{$n$-Vectors}

\subsection{Linear Dependence of $n$-vectors}

\subsection{Length \& Inner Product of $n$-vectors}

\subsection{Outer/Cross Product of 3-vectors}

\section{Finite-Dimensional Vector Spaces}

\subsection{$n$-dimensional Vector Spaces Over $\R$ and $\C$}

\subsection{Linear Dependence \& Generators}

\subsection{Simultaneous Linear Equation, Gaussian elimination}

\subsection{Bases \& Dimensions}

\subsection{Subspaces}

\subsection{Inner Products; Schwarz Inequality}

\subsection{Orthogonal Bases; Orthogonal Complements; Projections}

\subsection{Dual Spaces}

\subsection{Lines, Hyperplanes, Coordinate systems}

\subsection{Distance; Polarization Identity}

\section{Linear Transformations and Matrices}

\subsection{Linear Transformations; Elementary Properties (Image/Nullspace)}

\subsection{Addition, Composition, Scalar Multiplication}

\subsection{Matrices \& Linear Transformations}

\subsection{Nonsingular Linear Transformations; Inverses}

\subsection{Changes of Bases in Vector Spaces}

\subsection{Similarity of Matrices}

\subsection{Special Types of Square Matrices (Symmetric, Orthogonal, Triangular etc.)}

\subsection{Elementary Matrices}

\subsection{Rank of a Matrix}

\section{Determinants}

\subsection{Definition \& Elementary Properties}

\subsection{Permutations \& Uniqueness}

\subsection{Minors; Cofactors; Evaluation of Determinants}

\subsection{Applications (Dependence of Vectors; Volume of a Parallelopiped; Linear Equations)}

\subsection{Determinants of Products \& Inverses}

\subsection{Determinants of Special Types of Matrices}

\section{Bilinear and Quadratic Forms}

\subsection{Bilinear Mappings \& Forms}

\subsection{Quadratic Forms; Polarization}

\subsection{Equivalence of Quadratic Forms; Congruence of Matrices}

\subsection{Geometric Applications}

\section{Eigenvalues and Eigenvectors}

\subsection{Definitions}

\subsection{Similarity \& Diagonal Matrices}

\subsection{Orthogonal Reduction of Symmetric Matrices}

\subsection{Eigenvalues of Special Types of Matrices (Unitary, Hermitian, etc.)}

\subsection{Normal Matrices \& Spectral Theorem}

\subsection{Singular Values}

%================================================
\part{Probability}

\section{Combinatorial Analysis}

\subsection{Inclusion-Exclusion Principle}

\subsection{Permutations}

\subsection{Combinations}

\subsection{Multinomial Coefficients}

\section{Axioms of Probability}

\subsection{Sample spaces and events}

\subsection{Axioms of Probability}

\subsection{The uniform model}

\subsection{Probability as a continuous set function}

\section{Conditional probability and independence}

\subsection{Conditional probabilities}

\subsection{Bayes’ theorem}

\subsection{Independent events}

\section{Random variables and their distributions}

\subsection{Random variables}

\subsection{Discrete random variables, Indicator functions}

\subsection{Expected value}

\subsection{Expectation of a function of a random variable}

\subsection{Variance}

\subsection{Distribution function}

\subsection{Probability mass function}

\subsection{Probability density function}

\section{Discrete univariate distributions}

\subsection{Bernoulli and binomial distributions}

\subsection{Geometric and negative binomial distributions}

\subsection{Poisson distributions}

\subsection{Hypergeometric distributions}

\section{Continuous univariate distributions}

\subsection{Uniform distributions}

\subsection{Normal distributions}

\subsection{Exponential distributions}

\subsection{Gamma distributions}

\subsection{Cauchy distributions}

\subsection{Change of variables}

\section{Multivariate distributions}

\subsection{Independent random variables}

\subsection{Convolution and sums of independent random variables}

\subsection{Conditional distributions}

\subsection{Change of variables}

\subsection{Multivariate normal distributions}

\section{Properties of expectation and related matters}

\subsection{Linearity of expectation}

\subsection{Covariance}

\subsection{Variance of a sum}

\subsection{Correlation}

\subsection{Conditional expectation and the law of total expectation}

\subsection{Conditional variance and the law of total variance}

\subsection{Conditional covariance and the law of total covariance}

\subsection{Conditional expectation and prediction}

\subsection{Moment generating functions}

\section{Inequalities and limit theorems}

\subsection{Chebychev’s and Markov’s inequalities}

\subsection{Cauchy-Schwarz inequality}

\subsection{Jensen’s inequality}

\subsection{Weak and strong laws of large numbers}

\subsection{Central limit theorem}

%================================================
\part{Real Analysis}

\section{Basic set theory}

\subsection{Finiteness, countability, and uncountability}

\subsection{Ordered fields}

\subsection{Cauchy sequences}

\subsection{Construction of the real numbers}

\subsection{Completeness of the real numbers}

\subsection{Infimum and supremum}

\subsection{Liminf and limsup of sequences}

\subsection{Limit points of sets, and sequences}

\subsection{Norms, inner products, and metrics}

\subsection{Open sets and closed sets}

\subsection{Closure and boundary of a set}

\subsection{Compact sets}

\subsection{Bolzano–Weierstrass theorem}

\subsection{Heine–Borel theorem in Euclidean space}

\subsection{Intersection of nested compact sets}

\section{Continuous mappings}

\subsection{Sums, products, compositions of continuous mappings}

\subsection{Path-connected and connected sets}

\subsection{Images of compact and connected sets}

\subsection{Extreme value theorem}

\subsection{Intermediate value theorem}

\subsection{Uniform continuity}

\section{Differentiability}

\subsection{Chain rule}

\subsection{Product rule}

\subsection{Mean value theorem}

\subsection{Taylor’s theorem}

\subsection{Derivative matrix}

\subsection{Continuity of differentiable mappings}

\subsection{Differentiable paths}

\subsection{Directional derivatives}

\subsection{Inverse function theorem}

\subsection{Implicit function theorem}

\subsection{Maxima and minima}

\subsection{Constrained extrema and Lagrange multipliers}

\section{Integration}

\subsection{Existence and elementary properties of the Riemann integral}

\subsection{Fundamental theorems of calculus}

\subsection{Differentiation of limits of integration}

\subsection{Interchange of order of integration (Fubini’s theorem)}

\subsection{Differentiation of integrals with respect to a parameter in the integrand}

\subsection{Change of variables theorem}

\subsection{Polar, spherical, and cylindrical coordinates}

\subsection{Improper integrals}

\section{Sequences and series}

\subsection{Numerical sequence and series}

\subsection{Pointwise and uniform convergence of sequence of functions}

\subsection{The Weierstrass $M$-test}

\subsection{Differentiation and integration of series}

\subsection{Elementary functions (exponential, trigonometric, hyperbolic)}

\subsection{Contraction mapping theorem}

\subsection{The Stone-Weierstrass theorem}

\subsection{Dirichlet and Abel tests}

\subsection{Power series and Cesàro and Abel summability}

\end{document}